% Comparação dos Modelos

\chapter{Comparação dos Modelos}

Os três modelos avaliados apresentam diferenças consideráveis tanto em desempenho
quanto em custo computacional. A Tabela~\ref{tab:model-comparison} apresenta as principais
medidas obtidas durante os experimentos.

\input{summary\_table.tex}

A comparação das métricas de validação deve ser feita com cautela. O FFDSM e o
TFNet reportam \textit{val\_map50}, enquanto o ResCapsTCAttnFire reporta
\textit{val\_f1}. Portanto, os valores dessas métricas não constituem uma
comparação direta de desempenho preditivo. Dentro das métricas de mesmo tipo, o
FFDSM atingiu um melhor valor de \textit{mAP50} de 0.0917, enquanto o
TFNet atingiu 0.0493. O ResCapsTCAttnFire atingiu F1 de 0.7931, mas esse
resultado deve ser interpretado separadamente devido à diferença de métrica.

\section{Eficiência computacional}

A Figura~\ref{fig:efficiency} apresenta a comparação entre os modelos em termos
de GMACs, memória de pico, tempo médio por época e número de parâmetros.

\begin{figure}[htbp]
\centering
\includegraphics[width=\textwidth]{efficiency.png}
\caption{Comparação da eficiência computacional dos modelos.}
\label{fig:efficiency}
\end{figure}

O ResCapsTCAttnFire apresenta a menor quantidade de operações, com 4.633 GMACs,
e também o menor número de parâmetros, com 2.22 milhões. O FFDSM possui 9.271
GMACs e 9.24 milhões de parâmetros, enquanto o TFNet possui 31.656 GMACs
e 30.45 milhões de parâmetros. Assim, em termos de tamanho da arquitetura e
quantidade de operações, o ResCapsTCAttnFire apresenta a menor complexidade entre
os três modelos.

O consumo de memória apresenta um comportamento diferente. O TFNet
apresenta o menor pico de memória, com 3236.4 MB, seguido pelo
ResCapsTCAttnFire, com 3522.4 MB, e pelo FFDSM, com 4831.9 MB. Dessa forma, a
menor quantidade de parâmetros e GMACs do ResCapsTCAttnFire não se traduz
diretamente em menor utilização de memória durante o treinamento.

Em relação ao tempo médio por época, o FFDSM apresenta o menor valor, com
280.3 s por época. O TFNet apresenta 793.0 s, enquanto o
ResCapsTCAttnFire apresenta 1892.6 s por época. Isso mostra que a quantidade de
GMACs e de parâmetros, isoladamente, não é suficiente para explicar o tempo real
de execução observado durante o treinamento.

\section{Tempo de treinamento}

O tempo total de treinamento também deve ser considerado juntamente com o número
de épocas executadas. O FFDSM foi treinado por 129 épocas e apresentou um tempo
total de 10.04 horas. O TFNet foi treinado por 200 épocas e levou
44.05 horas. O ResCapsTCAttnFire, por sua vez, foi treinado por apenas 20 épocas
e apresentou um tempo total de 10.51 horas.

Apesar de o ResCapsTCAttnFire possuir o maior tempo médio por época, seu
treinamento completo apresentou duração semelhante à do FFDSM devido ao número
consideravelmente menor de épocas. O TFNet apresentou o maior tempo total
de treinamento, associado tanto ao elevado tempo por época quanto ao maior número
de épocas executadas.

\section{Métrica de validação}

A evolução da métrica de validação reportada para cada modelo é apresentada na
Figura~\ref{fig:metrics}.

\begin{figure}[htbp]
\centering
\includegraphics[width=\textwidth]{metric\_curves.png}
\caption{Evolução da métrica de validação reportada por cada modelo.}
\label{fig:metrics}
\end{figure}

O FFDSM apresenta uma evolução bastante instável de \textit{val\_map50}, com
oscilações consideráveis ao longo do treinamento. Apesar disso, observa-se uma
tendência geral de crescimento, atingindo o melhor valor de 0.0917 na época 125.

O TFNet apresenta oscilações ainda mais pronunciadas. Seu melhor valor de
\textit{val\_map50} foi 0.0493 na época 40. Após esse ponto, a métrica não
apresenta uma melhora consistente até o final do treinamento.

O ResCapsTCAttnFire apresenta comportamento diferente dos outros dois modelos,
uma vez que sua métrica reportada é F1. A curva apresenta uma tendência geral de
crescimento, partindo de aproximadamente 0.66 e chegando a 0.7931 na época 19.
Embora existam algumas quedas intermediárias, a evolução geral da métrica é
positiva ao longo das 20 épocas.

\section{Evolução da perda de treinamento}

A evolução da perda durante o treinamento é apresentada na
Figura~\ref{fig:loss}.

\begin{figure}[htbp]
\centering
\includegraphics[width=\textwidth]{loss\_curves.png}
\caption{Evolução da perda de treinamento dos modelos.}
\label{fig:loss}
\end{figure}

Os três modelos apresentam redução da perda durante o treinamento. O FFDSM
apresenta uma queda acentuada nas primeiras épocas, seguida por uma redução mais
gradual, atingindo uma perda final de aproximadamente 0.2231.

O ResCapsTCAttnFire também apresenta uma redução rápida da perda, chegando a
aproximadamente 0.2901 ao final das 20 épocas. Já o TFNet apresenta uma
redução mais lenta e permanece em aproximadamente 0.5731 ao final das 200 épocas.

Entretanto, os valores absolutos da perda não devem ser utilizados isoladamente
para determinar qual arquitetura apresenta melhor desempenho, especialmente caso
as funções de perda ou os procedimentos de treinamento utilizados pelos modelos
não sejam idênticos.

\section{Discussão}

Os resultados mostram diferentes compromissos entre complexidade computacional,
tempo de treinamento e desempenho de validação.

O ResCapsTCAttnFire apresenta a menor quantidade de parâmetros e GMACs entre os
três modelos. Seu número de parâmetros é aproximadamente 2.22 milhões, contra
9.24 milhões do FFDSM e 30.45 milhões do TFNet. Entretanto, essa
vantagem não se traduz em menor tempo por época: o ResCapsTCAttnFire apresenta
1892.6 s por época, aproximadamente 6.8 vezes o tempo observado para o FFDSM.
Esse comportamento indica que a contagem de GMACs não captura necessariamente
todos os fatores responsáveis pelo tempo efetivo de execução de uma arquitetura.

O FFDSM apresenta um compromisso diferente. Ele possui o menor tempo médio por
época entre os três modelos e, entre os modelos que reportam \textit{mAP50},
atingiu o maior valor observado, 0.0917. Seu principal custo computacional é o
maior pico de memória, de 4831.9 MB.

O TFNet apresenta 30.45 milhões de parâmetros e 31.656 GMACs, sendo o
modelo com maior complexidade computacional segundo essas duas medidas. Também
apresenta o maior tempo total de treinamento, 44.05 horas. Apesar de possuir o
menor pico de memória dos três modelos, com 3236.4 MB, seu \textit{mAP50} máximo
foi 0.0493.

É importante destacar que os resultados de F1 e mAP50 não devem ser comparados
numericamente como se fossem a mesma métrica. Assim, os dados permitem uma
comparação direta de eficiência computacional, mas a comparação de desempenho
preditivo entre os três modelos é limitada pela utilização de métricas diferentes.

Para uma comparação definitiva de desempenho, seria necessário avaliar os três
modelos sob o mesmo conjunto de validação, utilizando a mesma métrica e,
idealmente, o mesmo protocolo experimental.

